\documentclass[10pt,a4paper]{amsart}
\usepackage[margin=19mm]{geometry}
\usepackage{amsmath,amssymb,mathtools}
\usepackage[scr=rsfs,cal=euler]{mathalfa}
\usepackage{xfrac}
\usepackage[unicode,hidelinks,bookmarks=false]{hyperref}
\newcommand{\ie}{i.e.\ }
\newcommand{\cf}{cf.\ }
\newcommand{\resp}{respectively\ }
\newcommand{\Subsection}{\subsection}
\providecommand{\nobreakdash}{\nobreak\hbox{-}}
\setlength{\emergencystretch}{2em}
\multlinegap0pt
\newcounter{commentcounter}
\newcommand{\katie}[1]{\stepcounter{commentcounter}\textbf{Comment \arabic{commentcounter} (by Katie)}: {\textcolor{green}{#1}}}
\usepackage{stmaryrd}
\usepackage{enumitem}
\usepackage{ifthen}
\usepackage{mathabx}
\usepackage{mathtools}
    \hypersetup{urlcolor=RoyalBlue, linkcolor=RoyalBlue,  citecolor=black}
\usepackage{setspace}
\usepackage{tikz-cd}
\usepackage{xfrac}
\usepackage{cleveref}
\usepackage{wrapfig}
\newtheorem{theorem}{Theorem}[section]
\newtheorem{lemma}[theorem]{Lemma}
\newtheorem{corollary}[theorem]{Corollary}
\newtheorem{proposition}[theorem]{Proposition}
\newtheorem{conjecture}[theorem]{Conjecture}
\newtheorem{claim}[theorem]{Claim}
\newtheorem{addendum}[theorem]{Addendum}
\newtheorem{assumption}[theorem]{Assumption}
\newtheorem{question}[theorem]{Question}
\newtheorem{thmx}{Theorem}
\newtheorem{corx}[thmx]{Corollary}
\newtheorem{propx}[thmx]{Proposition}
\renewcommand{\thethmx}{\Alph{thmx}}
\renewcommand{\thecorx}{\Alph{corx}}
\newtheorem{definition}[theorem]{Definition}
\newtheorem{remark}[theorem]{Remark}
\newtheorem{remarks}[theorem]{Remarks}
\newtheorem{construction}[theorem]{Construction}
\newtheorem{setup}[theorem]{Setup}
\newtheorem{example}[theorem]{Example}
\newtheorem{examples}[theorem]{Examples}
\newtheorem*{ConjSinger}{The Singer Conjecture}
    \newtheorem{duplicate}{}
\newcommand*{\claimproofname}{My proof}
\newenvironment{claimproof}[1][\claimproofname]{\begin{proof}[#1]\renewcommand*{\qedsymbol}{\(\blacksquare\)}}{\end{proof}}
\DeclareMathOperator*{\Aster}{\text{\LARGE{\textasteriskcentered}}}
\DeclareMathOperator{\aster}{\text{\LARGE{\textasteriskcentered}}}
\DeclareMathOperator{\Aut}{\mathrm{Aut}}
\DeclareMathOperator{\Out}{\mathrm{Out}}
\DeclareMathOperator{\Inn}{\mathrm{Inn}}
\DeclareMathOperator{\Ker}{\mathrm{Ker}}
\DeclareMathOperator{\coker}{\mathrm{Coker}}
\DeclareMathOperator{\Coker}{\mathrm{Coker}}
\DeclareMathOperator{\Hom}{\mathrm{Hom}}
\DeclareMathOperator{\Ext}{\mathrm{Ext}}
\DeclareMathOperator{\Tor}{\mathrm{Tor}}
\DeclareMathOperator{\tr}{\mathrm{tr}}
\DeclareMathOperator{\im}{\mathrm{im}}
\DeclareMathOperator{\Fix}{\mathrm{Fix}}
\DeclareMathOperator{\orb}{\mathrm{Orb}}
\DeclareMathOperator{\End}{\mathrm{End}}
\DeclareMathOperator{\Irr}{\mathrm{Irr}}
\DeclareMathOperator{\Comm}{\mathrm{Comm}}
\DeclareMathOperator{\Isom}{\mathrm{Isom}}
\DeclareMathOperator{\Min}{\mathrm{Min}}
\DeclareMathOperator{\Core}{\mathrm{Core}}
\DeclareMathOperator{\bigset}{Big}
\DeclareMathOperator{\cay}{Cay}
\DeclareMathOperator{\diam}{diam}
\DeclareMathOperator{\hull}{hull}
\DeclareMathOperator{\link}{Lk}
\DeclareMathOperator{\map}{Map}
\DeclareMathOperator{\sym}{Sym}
\DeclareMathOperator{\Sym}{Sym}
\DeclareMathOperator{\Alt}{Alt}
\DeclareMathOperator{\alt}{Alt}
\DeclareMathOperator{\lat}{\mathrm{Lat}}
\DeclareMathOperator{\dom}{\mathrm{dom}}
\newcommand{\cala}{{\mathcal{A}}}
\newcommand{\calb}{{\mathcal{B}}}
\newcommand{\calc}{{\mathcal{C}}}
\newcommand{\cald}{{\mathcal{D}}}
\newcommand{\cale}{{\mathcal{E}}}
\newcommand{\calf}{{\mathcal{F}}}
\newcommand{\calg}{{\mathcal{G}}}
\newcommand{\calh}{{\mathcal{H}}}
\newcommand{\cali}{{\mathcal{I}}}
\newcommand{\calj}{{\mathcal{J}}}
\newcommand{\calk}{{\mathcal{K}}}
\newcommand{\call}{{\mathcal{L}}}
\newcommand{\calm}{{\mathcal{M}}}
\newcommand{\caln}{{\mathcal{N}}}
\newcommand{\calo}{{\mathcal{O}}}
\newcommand{\calp}{{\mathcal{P}}}
\newcommand{\calq}{{\mathcal{Q}}}
\newcommand{\calr}{{\mathcal{R}}}
\newcommand{\cals}{{\mathcal{S}}}
\newcommand{\calt}{{\mathcal{T}}}
\newcommand{\calu}{{\mathcal{U}}}
\newcommand{\calv}{{\mathcal{V}}}
\newcommand{\calw}{{\mathcal{W}}}
\newcommand{\calx}{{\mathcal{X}}}
\newcommand{\caly}{{\mathcal{Y}}}
\newcommand{\calz}{{\mathcal{Z}}}
\newcommand*{\frakg}{\mathfrak{G}}
\newcommand*{\fraks}{\mathfrak{S}}
\newcommand*{\frakt}{\mathfrak{T}}
\newcommand{\ulE}{\underline{E}}
\newcommand{\ulEG}{\underline{E}\Gamma}
\newcommand{\Efin}{E_{\mathcal{FIN}})}
\newcommand{\EfinG}{E_{\mathcal{FIN}}\Gamma}
\newcommand{\grp}{\mathbf{Grp}}
\newcommand{\set}{\mathbf{Set}}
\newcommand{\gtop}{\mathbf{Top}_\Gamma}
\newcommand{\orbf}{\mathbf{Or}_\calf(\Gamma)}
\newcommand{\bfE}{\mathbf{E}}
\newcommand{\Top}{\mathbf{Top}}
\newcommand{\spectra}{\mathbf{Spectra}}
\newcommand{\ko}{\mathbf{ko}}
\newcommand{\KO}{\mathbf{KO}}
\newcommand{\TRV}{\mathcal{TRV}}
\newcommand{\FIN}{\mathcal{FIN}}
\newcommand{\VC}{\mathcal{VC}}
\newcommand{\ALL}{\mathcal{ALL}}
\newcommand{\hocolim}{{\rm hocolim}}
\newcommand{\colim}{{\rm colim}}
\newcommand{\ind}{{\rm ind}}
\newcommand{\res}{{\rm res}}
\newcommand{\coind}{{\rm coind}}
\newcommand{\GL}{\mathrm{GL}}
\newcommand{\PGL}{\mathrm{PGL}}
\newcommand{\PGammaL}{\mathrm{P\Gamma L}}
\newcommand{\SL}{\mathrm{SL}}
\newcommand{\PSL}{\mathrm{PSL}}
\newcommand{\GU}{\mathrm{GU}}
\newcommand{\PGU}{\mathrm{PGU}}
\newcommand{\UU}{\mathrm{U}}
\newcommand{\SU}{\mathrm{SU}}
\newcommand{\PSU}{\mathrm{PSU}}
\newcommand{\Sp}{\mathrm{Sp}}
\newcommand{\PSp}{\mathrm{PSp}}
\newcommand{\OO}{\mathrm{O}}
\newcommand{\SO}{\mathrm{SO}}
\newcommand{\POmega}{\mathrm{P}\Omega}
\newcommand{\PSO}{\mathrm{PSO}}
\newcommand{\PGO}{\mathrm{PGO}}
\newcommand{\AGL}{\mathrm{AGL}}
\newcommand{\ASL}{\mathrm{ASL}}
\newcommand{\He}{\mathrm{He}}
\newcommand{\PG}{\mathrm{PG}}
\newcommand{\PSLp}{\mathrm{PSL}_2(\ZZ[\frac{1}{p}])}
\newcommand{\SLp}{\mathrm{SL}_2(\ZZ[\frac{1}{p}])}
\newcommand{\SLm}{\mathrm{SL}_2(\ZZ[\frac{1}{m}])}
\newcommand{\GLp}{\mathrm{GL}_2(\ZZ[\frac{1}{p}])}
\newcommand{\LM}{\mathrm{LM}}
\newcommand{\HLM}{\mathrm{HLM}}
\newcommand{\CAT}{\mathrm{CAT}}
\newcommand*{\lhalf}[1]{\overleftarrow{#1}}
\newcommand*{\rhalf}[1]{\overrightarrow{#1}}
\newcommand*{\sgen}[1]{\langle#1\rangle}
\newcommand*{\nest}{\sqsubseteq}
\newcommand*{\pnest}{\sqsubset}
\newcommand*{\conest}{\sqsupset}
\newcommand*{\pconest}{\sqsupsetneq}
\newcommand*{\trans}{\pitchfork}
\DeclareMathOperator{\Ad}{\mathrm{Ad}}
\DeclareMathOperator{\id}{id}
\newcommand{\onto}{\twoheadrightarrow}
\def\iff{if and only if }
\newcommand{\Stn}{\mathrm{St}_n}
\newcommand{\floor}[1]{\lfloor #1 \rfloor}
\DeclareMathOperator{\TSS}{\mathrm{TSS}_{\max}}
\DeclareMathOperator{\cosym}{\mathrm{coSym}}
\newcommand{\EE}{\mathbb{E}}
\newcommand{\KH}{\mathbb{K}\mathbf{H}} % Hyperbolic space over K
\newcommand{\RH}{\mathbb{R}\mathbf{H}} % Real hyperbolic space
\newcommand{\CH}{\mathbb{C}\mathbf{H}} % Complex hyperbolic space
\newcommand{\HH}{\mathbb{H}\mathbf{H}} % Quaternion hyperbolic space
\newcommand{\OH}{\mathbb{O}\mathbf{H}^2} % Cayley hyperbolic space
\newcommand{\Ffour}{\mathrm{F}_4^{-20}} % The other rank one group
\newcommand{\KP}{\mathbb{K}\mathbf{P}} % Projective plane over K
\newcommand{\RP}{\mathbb{R}\mathbf{P}} % Real projective plane
\newcommand{\CP}{\mathbb{C}\mathbf{P}} % Complex projective plane
\newcommand{\OP}{\mathbb{O}\mathbf{P}^2} % Cayley projective plane
\newcommand{\Covol}{\mathrm{Covol}}
\newcommand{\Vol}{\mathrm{Vol}}
\newcommand{\rank}{\mathrm{rank}}
\newcommand{\gd}{\mathrm{gd}}
\newcommand{\cd}{\mathrm{cd}}
\newcommand{\vcd}{\mathrm{vcd}}
\newcommand{\hd}{\mathrm{hd}}
\newcommand{\vhd}{\mathrm{vhd}}
\newcommand{\betti}{b^{(2)}}
\DeclareMathOperator{\lcm}{\mathrm{lcm}}
\DeclareMathOperator{\Char}{\mathrm{Char}}
\newcommand{\flag}{{\rm Flag}}
\newcommand{\wtX}{\widetilde{X}}
\newcommand{\wtXj}{\widetilde{X_J}}
\newcommand{\ulG}{\underline{\Gamma}}
\newcommand{\MM}{\mathbf{M}}
\newcommand{\repr}{\calr_\RR}
\newcommand{\repc}{\calr_\CC}
\newcommand{\reph}{\calr_\HH}
\newcommand{\nov}[3]{{\widehat{#1 #2}^{#3}}}
\newcommand{\cgr}[2]{#1 \llbracket #2 \rrbracket} %completed group ring
\def\Z{\mathbb{Z}}
\newcommand{\NN}{\mathbb{N}}
\newcommand{\ZZ}{\mathbb{Z}}
\newcommand{\CC}{\mathbb{C}}
\newcommand{\RR}{\mathbb{R}}
\newcommand{\QQ}{\mathbb{Q}}
\newcommand{\FF}{\mathbb{F}}
\newcommand{\KK}{\mathbb{K}}
\newcommand{\ode}{\mathrm{d}}
\newcommand*{\pde}[3][]{\ensuremath{\frac{\partial^{#1} #2}{\partial #3}}}
\usepackage{tikz}
\usetikzlibrary{arrows,quotes}
\tikzstyle{blackNode}=[fill=black, draw=black, shape=circle]
  \newcommand{\ang}[1]{\left\langle #1 \right\rangle}
\usepackage{amsaddr}
\begin{document}
\title[Growth rates, stable subgroups, and regular languages]{Growth rates, stable subgroups, and regular languages}
\author{Kaitlin Ragosta}
\date{}
\maketitle
\noindent\emph{Source and licence.} Growth rates, stable subgroups, and regular languages. Version arXiv:2607.01872v1 (CC BY 4.0). This material is distributed under the Creative Commons Attribution 4.0 International licence, http://creativecommons.org/licenses/by/4.0/, as recorded in the arXiv OAI licence metadata for this exact version and shown on the abstract page https://arxiv.org/abs/2607.01872v1. Full rights notice: licenses/arxiv-2607.01872-CC-BY-4.0.txt.\par\smallskip
\noindent\textit{Editorial scope.} Counted unit: the original \texttt{theorem} environment \texttt{thm:DFW-regular-lang} at source lines 524--536 together with its complete proof at lines 563--579; the delivered document keeps source lines 374--580 so that every referenced label and every citation resolves inside it. Native editable source is the paper's own concatenated TeX; the AMS wrapper is editorial formatting only. No publisher class, upstream script or unrelated artwork is redistributed.
\par\smallskip\noindent\textit{Reference note.} This article ships no compiled bibliography (biblatex); the entries delivered here are composed from the author's own BibTeX (.bib) fields, with no field invented and no entry dropped.
\begin{thmx}\label{thm:regular-language}
        Let $G$ be a group with finite generating set $A$, and let $H \leq G$ be a subgroup. Let $L_H$ be the language of geodesic words representing elements of $H$. The subgroup $H$ is stable if and only if $L_H$ is a regular language all of whose elements are $M$-Morse for some Morse gauge $M$. Moreover, if $H$ is stable, there is a sublanguage $J_H \subseteq L_H$ such that $J_H$ is regular and each element of $H$ is represented by exactly one word in $J_H$. 
\end{thmx}
\begin{corx}\label{corx:rational-growth}
    Let $H$ be a stable subgroup of a group $G$ with finite generating set $A$, and let $f_{H,A}(n) = |B_n(e) \cap H|$. There exist polynomials $P(x)$ and $Q(x) \in \mathbb{Q}[x]$ such that 
    \begin{align*}
        \sum_{n=0}^{\infty} f_{H,A} \cdot x^n = \frac{P(x)}{Q(x)}\text{.}
    \end{align*}
\end{corx}
Theorem \ref{thm:regular-language} gives a regular language characterization of the geometric property of stability, providing the first step towards resolving \cite[Question~1]{CRSZ}, which asked whether a subgroup is stable if and only if the language of $(k,c)$ quasi-geodesic words that start at the identity and end in $H$ is regular whenever $k$ is rational. This is the case for hyperbolic groups \cite{Holt-Rees-qgeod-regular,Sam-Patrick-Davide-qgeod-regular}. 

\cite[Question~5]{CRSZ} asked whether the growth rate of an infinite-index stable subgroup is always smaller than that of the ambient group; this was previously known when the ambient group is Morse local-to-global or admits a path system with a constricting element\cite{DFW,CRSZ,Growth-rate-gap-all-Mltg,Abdul-Alex-injective-contracting,Legaspi-path-system-growth}. We show that the non-positive curvature assumptions on $G$ can be removed if $H$ satisfies a sufficiently strong combination theorem.

\begin{thmx}\label{thmx:growth-rate-gap}
    Let $G$ be a group with finite generating set $A$, and let $H$ be a stable subgroup. Suppose there is an infinite-order element $g \in G$ such that $\langle H , g \rangle \cong H * \langle g \rangle$. The growth rate of $H$ with respect to $A$ is strictly smaller than the growth rate of $G$ with respect to $A$.
\end{thmx}

%We apply the \emph{transverse loxodromic elements} of \cite{Abbott-Hull} to deduce that a number of interesting stable subgroups of acylindrically hyperbolic groups satisfy a growth rate gap.
From this, we deduce a number of interesting examples.
\begin{corx}\label{corx:transverse-examples}
    For any finite set $A$ of $G$, $\lambda_{H,A} < \lambda_{G, A}$ in the following cases:
    \begin{enumerate}
        \item $G = \Out(\mathbb{F}_n)$ for $n \geq 3$ and $H$ is a convex cocompact subgroup.
        \item $G$ is a Torelli group $\mathcal{T}_g$ for $g > 2$ and $H$ is a convex cocompact subgroup of $\mathrm{MCG}(S_g)$ with $H < \mathcal{T}_g$.
        \item G is a genus $g$ handlebody group $\mathcal{H}_g$ for $g > 1$ and $H$ is a convex cocompact subgroup of $\mathrm{MCG}(\partial V_g)$ with $H < \mathcal{H}_g$.
    \end{enumerate}
\end{corx}

Finally, \cite[Question~3]{CRSZ} asked whether the membership problem is decidable for stable subgroups of Morse local-to-global groups and if their Theorem E, analogous to our Theorem \ref{thm:regular-language} in the Morse local-to-global case, had implications for the complexity of such an algorithm. In the final section, we show that this is not the case, even under additional desirable conditions.

\begin{propx}\label{propx:counterexample}
    There is a finitely presented, relatively hyperbolic, and Morse local-to-global group $G$ such that $G$ has a stable subgroup with undecidable membership problem. 
\end{propx}

\subsection{Acknowledgments}
 I would like to thank Carolyn Abbott for her mentorship and Abdul Zalloum for asking whether the properties in Theorem \ref{thm:regular-language} and Corollary \ref{corx:rational-growth} are intrinsic properties of stable subgroups. I would also like to thank Jacob Russell, David Futer, and Joshua Perlmutter for helpful conversations. I am grateful to my group at the 2026 Junior Research Retreat on Separability and Morseness in Kopp, namely Francesco Fournier-Facio, Huaitao Gui, Rafaela Ioannou, Joshua Perlmutter, Talia Schlomovich, and Bratati Som, for discussions related to the final section, and to the organizers of the retreat. I would especially like to thank Abdul Zalloum and Francesco Fournier Facio for several key observations and suggestions related to Theorem \ref{thm:regular-language} and Proposition \ref{propx:counterexample} respectively.

\section{Background}
Throughout the paper, generating sets are assumed to be symmetric. We first recall relevant definitions, beginning with stable subgroups.

\begin{definition}\cite[Definition~2.22]{CRSZ}
    Let $G$ be a group with a finite generating set $A$, let $M$ be a Morse gauge, and let $k\geq 0$. A subgroup $H\leq G$ is \textit{$(M,k)$-stable} in $\text{Cay}(G,A)$, if for every $h\in H$, every geodesic in $\text{Cay}(G,A)$ from $e$ to $h$ is $M$-Morse and contained in the $k$-neighborhood of $H$. A subgroup $H\leq G$ is a \textit{stable subgroup} if for any choice of finite generating set $A$ for $G$, there exist $M$ and $k$ such that $H$ is $(M,k)$-stable in $\text{Cay}(G,A)$.
\end{definition}
While this is stated differently than the original definition of stability in \cite{Durham-Taylor-stability}, the two are equivalent \cite{Cordes-Hume-stability}. We also recall the definition of the growth rate of a group or subgroup with respect to a particular finite generating set. 

\begin{definition}
    Let $G$ be a group with finite generating set $A$, and let $H \leq G$. The \emph{growth function} of $H$ with respect to $A$ is the function $f_{H,A}: \mathbb{N} \rightarrow \mathbb{N}$ where $f_{H,A}(n) = B_{n,A}(e) \cap H$, where $B_{n,A}$ is the ball of radius $n$ around the identity in $\text{Cay}(G,A)$. The \emph{growth rate} of $H$ with respect to $A$ is
\begin{align*}
    \lambda_{H,A} := \limsup_{n \rightarrow \infty} \sqrt[n]{f_{H,A}(n)}\text{.}
\end{align*}
\end{definition}

 %When $A'$ is a proper subset of $A$, however, $B_{n,A'}(e) \subseteq B_{n,A}(e)$ for every $n$, so $\lambda_{H,A'} \leq \lambda_{H,A}$.
In general, both the growth rate and the growth function are highly dependent on the choice of generating set $A$. We now recall several facts about languages and finite-state automata.
\begin{definition}
    Let $A$ be a finite set and $A^{\star}$ be the free monoid over $A$. A \emph{word} in $A$ is an element $w \in A^{\star}$. If $w = a_1 a_2 \cdots a_n$, the elements $a_i$ are the \emph{letters} of $w$. A \emph{language with alphabet $A$} is a set of words in $A^{\star}$. The \emph{word length} of $w$ is the number of letters which appear in it.
\end{definition}

When $A$ is the generating set of a group $G$, a word $w \in A^{\star}$ can also be viewed as an element of $G$. We call this element $\overline{w}$. Every path in $\text{Cay}(G,A)$ produces a word in $A^{\star}$ by concatenating labels of edges in the order they appear along the path, and conversely, every word in $A^{\star}$ produces a path in $\text{Cay}(G,A)$ by starting at the identity and following the path dictated by the letters in the word from left to right. A word in $A^{\star}$ is \emph{geodesic} if the associated path is a geodesic in $\text{Cay}(G,A)$ or, equivalently, if $w$ is a minimal length representative of $\overline{w}$ with respect to the generating set $A$. 

\begin{definition}
    Let $A$ be a finite set. A \emph{finite state automaton} with alphabet $A$ is a tuple $\mathcal{G} = (\Gamma, A, Y, s_0)$ where:
    \begin{itemize}
        \item $\Gamma$ is a finite directed graph. The vertices of $\Gamma$ are the \emph{states} of $\mathcal{G}$.
        \item Each edge of $\Gamma$ is labeled by an element of the alphabet $A$.
        \item $Y$ is a subset of the vertices of $\Gamma$. The vertices of $Y$ are the \emph{accept states} of $\mathcal{G}$, while the vertices not in $Y$ are the \emph{reject states}.
        \item $s_0$ is a vertex of $\Gamma$. We call $s_0$ the \emph{initial state} of $\mathcal{G}$.
    \end{itemize}
\end{definition}

\begin{definition}
    A \emph{directed path} in a finite state automaton $\mathcal{G} = (\Gamma, A, Y, s_0)$ is a sequence of edges $e_1, \cdots, e_n$ of $\Gamma$ such that the terminal vertex of $e_i$ is the initial vertex of $e_{i+1}$ for each $i \in \{1, \cdots, n-1\}$. A word $w = a_1 \cdots a_n$ in $A$ is \emph{read} by a path $e_1, \cdots, e_n$ in $\mathcal{G}$ if the label of $e_i$ is $a_i$ for each $i \in \{1, \cdots, n \}$. The \emph{language accepted by} $\mathcal{G}$ is the set of words in $A$ that are read by paths in $\mathcal{G}$ which start at the initial state $s_0$ and end at an accept state of $\mathcal{G}$. A language is \emph{regular} if it is the accepted language of some finite state automaton. 
\end{definition}

There is a growth function associated to a language.
\begin{definition}
    Given a language $L$ with alphabet $A$, the \emph{growth function} of $L$ is the function $f_{L}(n) : \mathbb{N} \rightarrow \mathbb{N}$ where
    \begin{align*}
        f_L(n) = |\{w \in L | \ell(w) \leq n\}|\text{.}
    \end{align*}
    The \emph{growth rate} of $L$ is
    \begin{align*}
        \lambda_{L} := \limsup_{n \rightarrow \infty} \sqrt[n]{f_{L}(n)}\text{.}
    \end{align*}
\end{definition}

\begin{theorem}\cite[Theorem~7]{Cannon-regular-geodesics}\label{thm:polynomial-growth}
    If $L$ is a regular language, there exist polynomials $P(x), Q(x) \in \mathbb{Q}[x]$ such that
    \begin{align*}
        \sum_{n=0}^{\infty} f_L(n) \cdot x^n  = \frac{P(x)}{Q(x)}\text{.}
    \end{align*}
\end{theorem}

\begin{definition}
    Given a finite directed graph $\Gamma$ with vertices $v_1, \cdots, v_n$, the \emph{adjacency matrix} for $\Gamma$ is the $(n \times n)$ matrix whose $(i,j)$th entry is 1 if there exists a directed edge connecting vertex $v_i$ to vertex $v_j$ and 0 otherwise. The \emph{adjacency matrix for a finite state automaton} $\mathcal{G} = (\Gamma, A, Y, s_0)$ is the adjacency matrix for the directed graph $\Gamma$. For either a directed graph $\Gamma$ or a finite state automaton $\mathcal{G}$, the \emph{Perron--Frobenius eigenvalue} $\rho_{\Gamma}$ or $\rho_{\mathcal{G}}$ is the eigenvalue of the adjacency matrix with largest absolute value.
\end{definition}

\begin{definition}
    A finite state automaton is \emph{pruned} if every state is the vertex of some path from the initial state to an accept state.
\end{definition}
Any finite state automaton can be pruned without altering the accepted language by removing any state which does not appear along any path from the initial state to an accept state. The Perron--Frobenius eigenvalue can be used to compute the growth rate of a regular language.

\begin{theorem}\cite[Theorem~3.6]{DFW}\label{thm:DFW-3-6}
    Let $\mathcal{G}$ be a pruned finite state automaton that accepts the regular language $L$. We have $\rho_{\mathcal{G}} \geq 1$ and
    \begin{align*}
        \rho_{\mathcal{G}} = \lim_{n \rightarrow \infty} \sqrt[n]{f_L(n)} = \lambda_L\text{.}
    \end{align*}
\end{theorem}
When the directed graph associated to one finite state automaton is contained in another, one can relate their Perron--Frobenius eigenvalues.
\begin{lemma}\cite[Lemma~3.2 and Theorem~3.4]{DFW}
    Let $\Gamma$ be a directed graph, and let $\Gamma'$ be a proper subgraph of $\Gamma$. If for every pair of distinct vertices $v,w \in \Gamma$, there exists a directed path in $\Gamma$ from $v$ to $w$ and $w$ to $v$, then $\rho_{\Gamma'} < \rho_{\Gamma}$.
\end{lemma}

We are interested in languages and automata because of the following connection with growth rates of subgroups. Corollary \ref{cor:growth-by-automaton} follows from Theorem \ref{thm:DFW-3-6}.
\begin{definition}
    Let $G$ be a group with finite generating set $A$, and let $L$ be a regular language with alphabet $A$. The language $L$ is \emph{geodesic} if for all $w \in L$, $w$ is a geodesic word in $\text{Cay}(G,A)$. It \emph{bijects with a subgroup} $H \leq G$ if for each $w \in L$, $\overline{w} \in H$ and the map $L \rightarrow H$ given by $w \rightarrow \overline{w}$ is a bijection.
\end{definition}
\begin{corollary}\label{cor:growth-by-automaton}
    Let $G$ be a finitely generated group with finite generating set $A$. Suppose there is a regular geodesic language $L$ with alphabet $A$ that bijects with a subgroup $H \leq G$. If $\mathcal{G}$ is a pruned finite state automaton that accepts $L$, then
    \begin{align*}
        \lambda_{H,A} = \lambda_L = \rho_{\mathcal{G}}\text{.}
    \end{align*}
\end{corollary}

Lastly, we will need to prove a straightforward lemma. It was shown in \cite{CRSZ} that if $H'$ is a finite-index subgroup of $H$, then $\lambda_{H,A} = \lambda_{H',A}$. In our setting, it will sometimes be necessary to pass to quotients by finite normal subgroups rather than to finite-index subgroups. We first make the following observation.

\begin{lemma}
    Let $H$ be a stable subgroup of a group $G$, and let $N$ be a finite normal subgroup of $G$. The image of $H$ under the quotient map $q : G \rightarrow G/N$ is a stable subgroup of $G/N$.
\end{lemma} 
\begin{proof}
    The quotient map $q: G \rightarrow G/N$ is both a homomorphism and a quasi-isometry. Stability is a quasi-isometry invariant, so the former condition implies that $q(H)$ is a subgroup, and the latter implies that it is stable.
\end{proof}
We can now prove the desired lemma.

\begin{lemma}\label{lem:quotient-by-finite-normal}
    Let $G$ be a group with finite generating set $A$, and let $H \leq G$. Let $N$ be a finite normal subgroup of $G$, and let $q : G \rightarrow G/N$ be the quotient map. The growth rates $\lambda_{H,A}$ and $\lambda_{q(H), q(A)}$ are equal.
\end{lemma}
%By $q(A)$, we simply mean that each element of $A$ is replaced with its image under the quotient, an element of the group $H/N$. If any element of $A$ was contained in $N$, we remove it, and we also remove duplicates if there are $a$ and $a' \in A$ with $a N = a' N$. Equivalently, one can think of $q(A)$ as a choice of coset representative for each coset of the form $aN$ for $a \in A$; if two words $w$ and $w'$ are given by the same sequence of cosets but using different coset representatives, then $wN = w'N$ by normality of $N$. We highlight that 
We highlight that $H$ need not be a proper subgroup.
\begin{proof}
    Let $h \in H$. First, notice that in general, $d_{A}(1,h) \geq d_{q(A)}(1, h N)$. Since for $h$ in the alphabet $A$ also yields a representative for the coset $hN$ in the alphabet $q(A)$, although if an element of $A$ is contained in $N$, the length may shorten. Thus every element of $B_{n,A}(e) \cap H$ yields an element of $B_{n, q(A)}(e) \cap q(H)$. Similarly, if $d_{q(A)}(1, hN) = n$ for some $n$, then the word given by lifting each element of $q(A)$ in a word for $hN$ to the corresponding element of $A$ has the same length and represents an element of $hN$. If $K = \max_{n \in N} d_A (e,n)$, this implies $d_A (1, h) \leq n +K$. Thus every element of $B_{n , q(A)}(e) \cap q(H)$ yields at least one element of $B_{(n+K), A}(e) \cap H$.
    
    Finally, notice that in the quotient, each $h \in G$ is identified with $|N|$ other elements, any of which may be contained in $B_{n,A}(e) \cap H$. Thus in one direction, we obtain a bound that $|B_{n,A}(e) \cap H| \leq |N||B_{n, q(A)}(e) \cap q(H)|$ and in the other, we obtain $|B_{n+K , A} (e) \cap H | \geq |B_{n , q(A)} (e) \cap q(H)|$. Since $K$ and $|N|$ do not depend on $n$, this implies $\lambda_{H,A} = \lambda_{q(H), q(A)}$. 
\end{proof}


\section{Regular languages}

To prove Theorem \ref{thm:regular-language} and Corollary \ref{corx:rational-growth}, we will apply the following theorem of \cite{DFW} to construct a regular language of geodesic words representing elements of our stable subgroup $H$.


\begin{theorem}{\cite[Theorem~6.7]{DFW}}\label{thm:DFW-regular-lang}
        Let $G$ be a hyperbolic group acting properly and cocompactly on a graph $\Upsilon$. Fix a basepoint $b \in \Upsilon$. There is a regular language $J_G$ with the following properties:
        \begin{enumerate}
            \item $J_G \rightarrow G \cdot b$ is a surjection with fibers of cardinality exactly $\mid \text{Stab}_G(b)\mid$.
            \item The words in $J_b$, i.e., the preimage of $b$, have length 1 or 0, and correspond to paths of length 0.
            \item Every word of length $n$ in $J_G - J_b$ corresponds to a length $n$ geodesic in $\Upsilon$, starting at $b$.
            \item For every quasi-convex subgroup $H\subset G$, the sublanguage $J_H \subset J_G$ of words mapping to $Hb$ is regular.
            \item Let $\rho_H$ be the Perron--Frobenius eigenvalue of the transition matrix for any pruned automaton accepting $J_H$. The growth rate $\lambda_H(\Upsilon)$ satisfies 
            \begin{align*}
                \Lambda_H(\Upsilon) = \lim_{n \rightarrow \infty} \sqrt[n]{f_{H,\Upsilon}(n)} =  \lim_{n \rightarrow \infty} \sqrt[n]{f_{L_H}(n)} = \rho_H \text{.}
            \end{align*}
        \end{enumerate}
    \end{theorem}


Stable subgroups are always hyperbolic. Of course, $H$ acts geometrically on $\text{Cay}(H,S)$ for any finite generating set $S \subseteq H$, but the choice of graph impacts the alphabet of the regular language obtained. Since we require a regular language in the alphabet $A$, which need not be contained in $H$, we will need $H$ to act geometrically on a subgraph of $\text{Cay}(G,A)$.

\subsection{A geometric action of $H$}\label{subsec:geom-action-of-H}

Throughout the section, let $G$ be a group with finite generating set $A$, and let $H$ be an $(M,k)$-stable subgroup of $G$. Define $\Upsilon$ to be the union of all geodesics in $\text{Cay}(G,A)$ which start and end in $H$. 
\begin{lemma}
    $H$ acts freely and cocompactly on $\Upsilon \subseteq \Gamma(G,A)$.
\end{lemma}
\begin{proof}
    Let $x \in \text{Cay}(G,A)$ lie on a geodesic $\gamma$ from $p \in H$ to $q \in H$. The entire group $G$ acts by isometries, so for any $h \in H$, $h \cdot \gamma$ is a geodesic from $h \cdot p \in H$ to $h \cdot q \in H$ containing $h \cdot x$. Thus left-multiplication gives a group action $H \curvearrowright \Upsilon$. That this action is free follows immediately from the fact that $G \curvearrowright \text{Cay}(G,A)$ is free.
    
    Let $K = \Upsilon \cap B_{k+1, A}(e)$. Since $B_{k+1, A}(e)$ contains finitely many vertices and edges, so does $K$. Let $x_1$ and $x_2$ be two vertices connected by an edge $s$ all of which lie on a geodesic $\gamma \subset \Upsilon$ as before. Since $H$ is stable, there is some $h \in H$ such that $d_A (h, x_1) \leq k$. Since $H$ acts by isometries, $d_A (e, h^{-1} \cdot x) \leq k$. By construction, $d_A (e , h^{-1} x_2) \leq k+1$. The edge between the $h^{-1} \cdot x_i$, namely $h^{-1} \cdot s$, lies on the geodesic $h^{-1} \cdot \gamma$, so it is in $K$. Thus $\Upsilon = H \cdot K$ for a compact set $K$.
\end{proof}

\subsection{A regular language from the action}\label{subsec:regular-language-alphabet}
    We must now verify that the regular language obtained from this action has the correct alphabet. We also make explicit a detail which does not appear in \cite[Theorem~6.7]{DFW} but does follow from their method of proof.
    \begin{lemma}\label{lem:right-alphabet}
        In the case of $H$ acting on $\Upsilon$ as defined in the previous section, Theorem \ref{thm:DFW-regular-lang} gives a regular language with alphabet $A$ which bijects $H$. Furthermore, the language of geodesic words in $A$ which represent elements of $H$ is regular.
    \end{lemma}

    Before we begin, we recall the following theorem of Gersten and Short.
    \begin{theorem}\cite[Theorem~2.2]{Gersten-Short}
        Let $G$ be a hyperbolic group with an automatic structure $L_G$, and let $H$ be a quasi-convex subgroup. The sublanguage $L_H \subset G$ of words which map $H$ to itself is a regular language.
    \end{theorem}
    We can now prove the lemma.

    \begin{proof}
        Essentially, we proceed by carefully examining the proof in \cite{DFW} of Theorem \ref{thm:DFW-regular-lang}. First, we note that Construction 6.1 of \cite{DFW} is not necessary. Construction 6.1 replaces $\Upsilon$ with a graph $\Upsilon^{*}$ on which the group acts freely, but in our setting, $H$ already acts freely on $\Upsilon$. Thus we can begin with Construction 6.4, which replaces the group $H$ with a group $H^{+}$ which acts transitively on a graph $\Upsilon^{+}$. The subgroup $H$ does \emph{not} act transitively on $\Upsilon$ in our setting; some vertices in $\Upsilon$ are elements of $H$ while some are not, and these cannot lie in the same $H$-orbit.

        \textbf{Construction 6.4:} Attach 2-cells to $\Upsilon$ as follows. Choose a single representative from each $H$-orbit of based cycles, and attach a 2-cell along it. Extend this $H$-equivariantly to obtain a simply connected 2-complex with a free $H$-action. Call this $\Upsilon^{*}$. Let $D$ be the quotient of $H \setminus \Upsilon^{*}$ obtained by identifying all 0-cells. Define $H^{+} := \pi_1(D) \cong H * F_r$, where $F_r$ is a free group whose generators are in one-to-one correspondence with edges in a spanning tree for $H \setminus \Upsilon^{*}$. These come from orbits of edges in $\Upsilon$ which are not contained in $H$. Consider $\Tilde{D}$, the universal cover of $D$, which is a tree of copies of $\Upsilon^{*}$. Let $\Upsilon^{+}$ be the 1-skeleton of $\Tilde{D}$. The group $H^{+}$ is a deck group, which acts transitively on the vertices of $\Upsilon^{+}$.

        We can now proceed to Construction 6.5, which builds the regular language. 

        \textbf{Construction 6.5:} Choose a generating set $S^{+}$ for $H^{+}$ as follows: each generator corresponds to a closed path in the $1$-skeleton of $D$ consisting of a single edge. It is clear from the construction that $S^{+}$ generates $H^{+}$; $D$ has only one vertex, so every path in the fundamental group is a concatenation of the loops in $S^{+}$.
        
        This differs very slightly from the proof in \cite{DFW}; there, the generators are of the form $y e y'$ where $e$ is an edge in our sense and $y$ and $y'$ are a separate class of edges called \emph{tiny edges}. The tiny edges are added in Construction 6.1 to create a free action, and we did not need to implement that construction, so our complex $D$ does not have any tiny edges. The set $S^{+}$ is symmetric by construction.

        Recall that $D$ is a quotient of $\Upsilon^{*}$ where all 0-cells are identified, and $\Upsilon^{*}$ has $\Upsilon$ as its 1-skeleton. In particular, each edge in $D$ was an edge in $\Upsilon$, so each element of $S^{+}$ was a single element of $A$. The set $S^{+}$ is finite because the action of $H$ on $\Upsilon$ was cocompact.

        Since $H^{+} \cong H * F_r$ is hyperbolic, \cite[Theorem~11.27]{Cannon-regular-language} provides a geodesic regular language $L^{+}$ which maps bijectively to $H^{+}$. The alphabet of this language is the generating set $S^{+}$. Let $J_H$ be the sublanguage mapping bijectively to $H \subseteq H^{+}$. Since $H$ is quasi-convex in $H^{+} \cong H * F_r$, the sublanguage $J_H$ is regular by \cite[Theorem2.2]{Gersten-Short}.

        Finally, we note that the language of all geodesics with respect to a given finite generating set in a hyperbolic group is regular \cite{Cannon-regular-geodesics}. Again, $H$ is quasi-convex in the hyperbolic group $H^{+}$, so the sublanguage of geodesic words in the alphabet $S^{+} \subseteq A$ which represent elements of $H$ is also regular by \cite[Theorem~2.2]{Gersten-Short}. By construction, $S^{+}$ includes every letter of $A$ which appears on any geodesic between elements of $H$, so this is precisely the language of geodesic words in $A$ representing elements of $H$.
 \end{proof}


\begin{thebibliography}{99}
\bibitem[]{Abdul-Alex-injective-contracting} Sisto, Alessandro; Zalloum, Abdul. Morse subsets of injective spaces are strongly contracting. Groups, Geometry and Dynamics 19 (2025) 1425--1443
\bibitem[]{CRSZ} Cordes, Matthew; Russell, Jacob; Spriano, Davide; Zalloum, Abdul. Regularity of {M}orse geodesics and growth of stable subgroups. Journal of Topology 15 (2022) 1217–1247
\bibitem[]{Cannon-regular-geodesics} Cannon, James W.. The combinatorial structure of cocompact discrete hyperbolic groups. Geometriae Dedicata 16 (1984)
\bibitem[]{Cannon-regular-language} Cannon, J.W.. The theory of negatively curved spaces and groups. Ergodic theory, symbolic dynamics and hyperbolic spaces (1991) 315--369
\bibitem[]{Cordes-Hume-stability} Cordes, Matthew; Hume, David. Stability and the {M}orse boundary. Journal of the London Mathematical Society 95 (2017) 963--988
\bibitem[]{DFW} Dahmani, Fran\c{c}ois; Futer, David; Wise, Daniel. Growth of quasiconvex subgroups. Mathematical Proceedings of the Cambridge Philosophical Society 163 (2019) 505-530
\bibitem[]{Durham-Taylor-stability} Durham, Matthew Gentry; Taylor, Samuel J. Convex cocompactness and stability in mapping class groups. Algebraic \& Geometric Topology 15 (2015) 2837--2857
\bibitem[]{Gersten-Short} Gersten, S.M.; Short, H.B.. Rational subgroups of biautomatic groups. Annals of Mathematics 134 (1991) 125--158
\bibitem[]{Growth-rate-gap-all-Mltg} Han, Suzhen; Liu, Qing. Growth rate gap for stable subgroups. arXiv:2412.11244 (2024)
\bibitem[]{Holt-Rees-qgeod-regular} Holt, Derek; Rees, Sarah. Regularity of quasi-geodesics in a hyperbolic group. International Journal of Algebra and Computation 13 (2003) 585--596
\bibitem[]{Legaspi-path-system-growth} Legaspi, Xabier. Growth of quasi-convex subgroups in groups with a constricting element. Groups, Geometry and Dynamics 18 (2024) 1469--1505
\bibitem[]{Sam-Patrick-Davide-qgeod-regular} Hughes, Sam; Nairne, Patrick; Spriano, Davide. Regularity of quasi-geodesics characterizes hyperbolicity. Proceedings of the Royal Society of Edinburgh Section A: Mathematics (2025) 1--14
\end{thebibliography}
\end{document}
